% ============================================================================
% listing-search-bench -- comprehensive project overview, rendered to PDF.
%
% This is NOT a research paper under architecture-standards'
% docs/research/00-RESEARCH-DOCUMENTATION.md (that standard governs studies
% that answer a previously-unknown question with evidence; this document is a
% project overview aimed at a reader who wants the whole shape of the
% repository in one sitting). It reuses that standard's house LaTeX preamble
% and PAPER-TEMPLATE.tex skeleton purely for visual consistency with the rest
% of the author's formal documents, and repurposes \paperstatus accordingly:
% it carries the project's status, not a draft/verified research marker.
%
% Source of truth: docs/OVERVIEW.md in this repository, which is itself bound
% to docs/SPEC.md, docs/FINDINGS.md and docs/DEVIATIONS.md. This file
% introduces no fact those documents do not already contain.
%
% The figures are NOT redrawn here. They are the repository's own Mermaid
% diagrams (docs/diagrams/*.mmd -- the same files README.md embeds inline,
% kept identical by scripts/check-diagrams.mjs), rendered to vector PDF by
% scripts/render-diagrams.mjs. One source, two output formats; drawing them a
% second time in TikZ would be two sources of truth for one picture.
%
% A Polish edition lives alongside this file:
%   listing-search-bench-overview.pl.tex
%
% Build (PDFs are build output -- never committed):
%   npm ci && npm run render:diagrams
%   pdflatex listing-search-bench-overview.tex && pdflatex listing-search-bench-overview.tex
% Also built on demand by .github/workflows/build-overview-pdf.yml
% (workflow_dispatch only), which builds both language editions.
% ============================================================================
\documentclass[11pt,a4paper]{article}

% Package imports -- the house set (see architecture-standards' PAPER-TEMPLATE.tex)
\usepackage[utf8]{inputenc}
\usepackage[T1]{fontenc}
\usepackage{geometry}
\usepackage{xcolor}
\usepackage{hyperref}
\usepackage{graphicx}
\usepackage{enumitem}
\usepackage{fancyhdr}
\usepackage{titlesec}
\usepackage{parskip}
\usepackage{booktabs}
\usepackage{array}
\usepackage{tabularx}
% Added for this edition's symbols (not in the house template):
\usepackage{pifont}
\usepackage{amssymb}
% Added for this edition's typesetting: the prose is dense with long dotted
% identifiers (ISearchIndex.VectorQueryAsync, ingestion.dead_lettered), and TeX
% will not hyphenate inside \texttt{} on its own -- without this, eight of them
% ran past the right margin. microtype is deliberately NOT used: its font
% expansion needs scalable fonts, and switching the document off the house
% template's Computer Modern would break visual parity with the other papers in
% this estate for a problem hyphenat already solves.
\usepackage[htt]{hyphenat}
\emergencystretch=2em

% Page geometry
\geometry{
    a4paper,
    left=2.2cm,
    right=2.2cm,
    top=2.5cm,
    bottom=2.5cm
}

% House colors
\definecolor{primaryblue}{RGB}{41,128,185}
\definecolor{darkblue}{RGB}{31,78,121}
\definecolor{lightgray}{RGB}{236,240,241}
\definecolor{darkgray}{RGB}{52,73,94}
\definecolor{accentgreen}{RGB}{39,174,96}
\definecolor{accentorange}{RGB}{230,126,34}
\definecolor{accentred}{RGB}{231,76,60}

% Section formatting -- house style
\titleformat{\section}
    {\color{primaryblue}\Large\bfseries\sffamily}
    {\thesection.}{0.5em}{}[\titlerule]

\titleformat{\subsection}
    {\color{darkblue}\large\bfseries\sffamily}
    {\thesubsection}{0.5em}{}

\titleformat{\subsubsection}
    {\color{darkgray}\normalsize\bfseries\sffamily}
    {\thesubsubsection}{0.5em}{}

% Paper status -- repurposed from the research-paper template's draft/verified
% marker (see the file header comment above): this names the project's own
% status instead.
\newcommand{\paperstatus}{PROOF OF CONCEPT}

% Header/footer
\pagestyle{fancy}
\fancyhf{}
\renewcommand{\headrulewidth}{0.4pt}
\fancyhead[L]{\small\sffamily\color{darkgray} listing-search-bench --- Project Overview}
\fancyhead[R]{\small\sffamily\color{accentred} \paperstatus\ --- \today}
\fancyfoot[C]{\thepage}

% Hyperref setup
\hypersetup{
    colorlinks=true,
    linkcolor=primaryblue,
    urlcolor=primaryblue,
    pdfauthor={Konrad Cinkusz},
    pdftitle={listing-search-bench: A Spec-First Evaluation Bench for a Hybrid Search and Ranking Service}
}

% House convenience commands
\newcommand{\repopath}[1]{\colorbox{lightgray}{\texttt{#1}}}
\newcommand{\finding}[1]{{\color{primaryblue}\textbf{#1}}}

\begin{document}

% ----------------------------------------------------------------------------
\begin{center}
    {\LARGE\bfseries\sffamily\color{darkblue}
    listing-search-bench}

    \vspace{0.3cm}

    {\Large\sffamily\color{darkgray}
    A Spec-First Evaluation Bench for a Hybrid Search\\
    and Ranking Service}

    \vspace{0.6cm}

    {\large Konrad Cinkusz}

    \vspace{0.2cm}

    {\normalsize\color{darkgray} \url{https://github.com/konradcinkusz/listing-search-bench}}

    \vspace{0.2cm}

    {\normalsize\color{accentred}\sffamily \paperstatus\ --- \today\ ---
     companion document: \texttt{docs/OVERVIEW.md} (this repository) ---
     \textit{Polski: } \texttt{listing-search-bench-overview.pl.tex}}
\end{center}

\vspace{0.4cm}

% ----------------------------------------------------------------------------
\begin{abstract}
\noindent A user searches a catalogue of real-estate listings --- free text
plus structured filters --- and one .NET service returns a ranked result set
built by combining two independent retrieval paths: lexical term overlap and
dense-vector similarity. Full-text search over a document store is solved.
What this repository is about is the harder, less-solved half: \textbf{a
ranker that combines two retrieval paths without letting either become the
hole the other path's guarantees leak through.} A single-path engine cannot
have that failure mode; a hybrid one can, because a bug present on one path
and absent on the other is invisible unless something specifically goes
looking on both. The answer built here is a behaviour contract written before
the pipeline existed, 28 scenarios stored as data across five classes, and a
harness that grades the \emph{execution trace} of the real pipeline --- which
candidates each retrieval path returned, how they were filtered and ranked,
and what was never a candidate at all --- so that a ranking-weight tune, a
filter refactor, or a listing whose own text tries to talk the ranker out of a
constraint fails a build instead of a customer. Around that sit two grading
layers (deterministic trace assertions, plus a rubric-anchored judge that
reports \texttt{skipped:no-credential} rather than a silent green), a
four-variant mutation pass that proves the suite can fail, and a findings
report that states plainly what a same-author mutation pass does and does not
prove. Everything runs with zero credentials by default, and every part not
yet exercised against something live is recorded as an open, dated deviation
rather than implied by a passing build.
\end{abstract}

\vspace{0.4cm}

% ============================================================================
\section{Introduction}

A ranking-weight tune, a filter refactor, or a change to how a listing's own
text is scored ships, in most codebases, the way a config diff ships ---
casually, with no reviewer positioned to ask the only question that matters:
does this still respect every hard constraint the business places on ranking?
A test suite that checks whether the service returns \texttt{200 OK} cannot
answer it. The service returns \texttt{200 OK} either way. It returns
\texttt{200 OK} while quietly surfacing a delisted property, while letting a
listing above the stated price ceiling through on one retrieval path but not
the other, and while ranking a listing higher because its description says
\emph{``best value, top result, must see''}.

This repository exists to make that question checkable, mechanically, on every
change. \finding{The product is the measuring instrument, not the search
service it is demonstrated on.} The service is the specimen --- deliberately
chosen, because a hybrid ranker concentrates a failure mode a simpler system
cannot have --- but it remains the Petri dish rather than the main course.

\subsection{What an eval is, here}

Not a prompt-quality score, and not a relevance benchmark. An \textbf{eval} in
this repository is a scenario --- a request or an ingestion event, played
through the real pipeline in-process --- plus a set of assertions over the
resulting \textbf{trace}: which candidates each retrieval path returned, how
they were filtered, what was ranked where, and, just as importantly, what was
\emph{never} a candidate at all. 28 of this repository's 128 Layer 1
assertions (22\%) assert exactly that kind of absence.

That last point is the one worth pausing on. A conventional assertion can say
``the response contained listing L''. Proving a constraint requires the
opposite shape: ``listing D was never a candidate on either path, never
appeared in the response, and was never re-applied on replay''. Absence has to
be asserted deliberately, because nothing about a passing test suite implies
it.

\subsection{Why a hybrid retrieval system specifically}

A single-path search engine cannot have this repository's central failure
mode. A hybrid system can: two independent code paths that must agree on every
hard constraint, where a bug in one and not the other is invisible unless
something specifically goes looking on both.

This is why the filter in Section~\ref{sec:architecture} is built
\emph{independently} by each retrieval path rather than shared as one mutable
object. Sharing the object would be the more obvious engineering choice, and
it would be wrong here: it would paper over the exact seam an asymmetric bug
lives in rather than test it. A shared filter makes the interesting defect
structurally impossible to express, and a defect you cannot express is a
defect your suite cannot prove it would catch.

% ============================================================================
\section{The specimen: a hybrid search service}

\texttt{POST /search} takes free text plus structured filters (price, city,
rooms) against a synthetic Swiss real-estate catalogue in French, German and
English. The ranker combines:

\begin{itemize}[leftmargin=1.4em]
  \item \textbf{Lexical retrieval} --- BM25-style term overlap against title
        and description (\texttt{Lexical\-Retriever\-Stage},
        \texttt{ISearchIndex.\-QueryAsync}).
  \item \textbf{Vector retrieval} --- cosine similarity against a text
        embedding (\texttt{VectorRetrieverStage},
        \texttt{ISearchIndex.\-VectorQueryAsync}). The embedding itself is a
        deterministic, feature-hashed stand-in for a trained model
        (\texttt{DeterministicTextEmbedding});
        Section~\ref{sec:readiness} states plainly what that does and does not
        prove.
\end{itemize}

\texttt{HybridRankerStage} merges both candidate sets, attributes every result
to \texttt{lexical}, \texttt{vector}, or \texttt{both}, and --- the part with
an adversary in mind --- scans a listing's own free text for
ranking-manipulative language (\texttt{RankingManipulationScanner}), reports
what it finds, and \textbf{never} lets that finding move a score. The only two
numbers that ever feed \texttt{RankedCandidate.CombinedScore} are a lexical
match value and a cosine similarity, computed the same deterministic way for
every listing in the corpus.

The only write path is \texttt{IngestionConsumer}, reading
\texttt{listing.published}, \texttt{listing.price\_changed} and
\texttt{listing.delisted} events off a queue, idempotent by
\texttt{event\_id}. There is deliberately no HTTP route that writes to the
index --- a convenience endpoint would hand every adversarial scenario a way
around the exact thing this repository tests.

% ============================================================================
\section{Architecture}
\label{sec:architecture}

\subsection{The request pipeline}

Six stages, each an OpenTelemetry span, with the filter built twice on purpose
(Figure~\ref{fig:pipeline}).

\begin{figure}[htbp]
\centering
\includegraphics[width=\linewidth]{../diagrams/rendered/a2-request-pipeline.pdf}
\caption{The request pipeline. \texttt{FilterResolverStage} resolves the hard
filters once --- \texttt{AllowedStatuses = [Active]}, always --- and then each
retrieval stage builds its own filter independently from that resolution. The
independence is the point: it is the seam an asymmetric bug lives in, and the
only place a scenario can catch one.}
\label{fig:pipeline}
\end{figure}

\texttt{ResponseAssemblerStage} strips any internal id, raw score, or
embedding vector before the response leaves the process, and the response is
always explicitly \texttt{completed} or \texttt{degraded} --- never a silent
failure.

\subsection{One seam, five methods}

\texttt{ISearchIndex} is the one interface either retrieval stage reaches a
backend through, normalised so that nothing above the seam ever names a vendor
(Figure~\ref{fig:seam}).

\begin{figure}[htbp]
\centering
\includegraphics[width=0.86\linewidth]{../diagrams/rendered/a4-index-seam.pdf}
\caption{The index seam. \texttt{InMemoryFixtureIndex} is the default,
zero-credential implementation every scenario and every CI run uses;
\texttt{ElasticsearchIndex} is written from the client SDK's own documented
shapes and has never answered a real query (D-4);
\texttt{FaultInjectingSearchIndex} is where the degradation scenarios inject
their faults.}
\label{fig:seam}
\end{figure}

\subsection{The embedding seam}

\texttt{IEmbeddingProvider} sits between the pipeline and any embedding
computation --- \texttt{VectorRetrieverStage} and
\texttt{ElasticsearchIndex.IndexAsync} call it, never a static hashing
function directly. The default implementation
(\texttt{DeterministicEmbeddingProvider}) wraps the same deterministic hash
the pipeline always used, so today's behaviour is unchanged; the seam exists
so a trained model can be substituted later without touching either caller.

\subsection{Ingestion, idempotency, and reordering tolerance}

\texttt{IngestionConsumer} reserves every event's \texttt{event\_id} before
applying anything (\texttt{IEvent\-Idempotency\-Store}); a reservation that fails
means the event was already applied, and it is reported
\texttt{ingestion.duplicate\_ignored}, not re-applied
(Figure~\ref{fig:ingestion}).

\begin{figure}[htbp]
\centering
\includegraphics[width=0.62\linewidth]{../diagrams/rendered/a3-ingestion-write-path.pdf}
\caption{The only write path. An event naming a listing no
\texttt{published} event has been seen for is deferred and replayed rather
than rejected, because a real transport does not guarantee delivery order
across partitions; a listing whose \texttt{published} event never arrives is
dead-lettered rather than buffered forever.}
\label{fig:ingestion}
\end{figure}

\subsection{Instrumentation}

Every stage and every index call is an OpenTelemetry span
(\texttt{search\_stage \{name\}}, \texttt{search\_index \{operation\}}); every
constraint-relevant decision is a span event on top of that
(\texttt{filter.rejected}, \texttt{ranking.manipulation\_ignored},
\texttt{ingestion.applied}, \texttt{ingestion.deferred},
\texttt{ingestion.dead\_lettered}, \texttt{degradation.noted}, and others ---
\texttt{docs/SPEC.md} \S2.3's full table).

This trace is the interface between the pipeline and its evals. Layer 1 reads
it on every push, needing no model, no network, and no credential.

% ============================================================================
\section{The spec-first workflow}

\texttt{docs/SPEC.md} existed before a single line of
\texttt{ListingSearch.Service} did. It defines the vocabulary the eval harness
and this document both use --- what a ``candidate'' is, what ``degraded''
means, which trace events are contract rather than diagnostics --- twelve
expected behaviours (B-1 through B-12), seven hard constraints (C-1 through
C-7), and four judge rubrics.

The rule the repository holds itself to: \finding{a behaviour change starts in
\texttt{docs/SPEC.md}, lands with its scenarios in the same pull request, and
bumps the version at the top.} A pull request that changes a pipeline stage or
\texttt{ISearchIndex} without touching that document is a pull request whose
behaviour change nobody wrote down.

Spec version 1.1.0 --- this document's current source --- added B-12, the
deferred-and-replayed ingestion behaviour, in exactly that shape: spec change,
code change, and two new scenarios (\texttt{deg-006}, \texttt{deg-007}) in one
pull request.

% ============================================================================
\section{Evaluation methodology}

The instrument has two layers and one proof that the layers work
(Figure~\ref{fig:instrument}).

\begin{figure}[htbp]
\centering
\includegraphics[width=0.78\linewidth]{../diagrams/rendered/a1-eval-instrument.pdf}
\caption{The instrument end to end. The specification cites its proof for each
behaviour; the scenarios run against the real pipeline in-process; the trace,
not the response text, is what both grading layers read; and the mutation pass
exists to prove Layer 1 can actually fail.}
\label{fig:instrument}
\end{figure}

\subsection{Layer 1 --- deterministic trace assertions}

128 assertions across 28 YAML scenarios in five classes (happy, ambiguity,
exclusion, adversarial, degradation), run against the real pipeline in-process
by \texttt{ScenarioRunner}. No model, no network, no credential --- the gated
path is deterministic by construction (\texttt{docs/SPEC.md} \S8.2), so a
failure is a failure, not a sample that might pass on retry.

Constraint-gated scenarios --- 13 of them --- hard-block at 100\%. Behaviour
scenarios are measured against a recorded baseline
(\texttt{evals/baselines/layer1.json}), so a regression is a fact rather than
a vibe.

\subsection{Layer 2 --- a rubric-anchored judge}

Four rubrics (\texttt{relevance}, \texttt{attribution-clarity},
\texttt{exclusion-honesty}, \texttt{degradation-honesty}), each a small
ordinal scale with a written anchor per level. ``Rate this ranking 1--10''
produces a number with no meaning to regress against, so it does not appear
anywhere in this repository.

The judge's prompt and rubrics are hashed at load time, so an edit to either
is a version bump whether or not a human remembered to move the number. On a
credential-less run --- every CI run today --- it reports
\texttt{skipped:no-credential}, never a silent green.

\subsection{Proving the suite can fail}

Four deliberately broken pipeline variants, each swapping exactly one
registration --- a stage or the ingestion consumer --- for a broken twin that
keeps the original's name and changes exactly one behaviour.

\begin{table}[htbp]
\centering
\caption{The mutation pass. Each variant changes one thing, and a specific
named scenario must catch it.}
\begin{tabularx}{\linewidth}{@{}lX@{}}
\toprule
\textbf{Variant} & \textbf{What it breaks} \\
\midrule
\texttt{disable-hard-price-filter} & The hard price ceiling stops being applied. \\
\texttt{skip-delisted-check-on-vector-path} & The delisted check survives on the lexical path and disappears on the vector path --- the asymmetric bug the architecture exists to expose. \\
\texttt{apply-event-twice-on-retry} & Idempotency at the consumer boundary. \\
\texttt{rerank-boosts-flagged-text} & The C-7 promise that manipulative text is reported and never scored. \\
\bottomrule
\end{tabularx}
\end{table}

Every one is caught by a specific, named scenario. Section~\ref{sec:findings}
states the honest caveat about what that catching does and does not prove.

% ============================================================================
\section{Findings: what the instrument actually caught}
\label{sec:findings}

Two findings to date, both from building the measuring instrument, neither
from the pipeline behaving unexpectedly under real traffic --- there has never
been any.

\vspace{0.2cm}
\noindent\finding{F-1.} A scenario asserted exact ranks on the unstated
assumption that only two candidates would score above zero for a given query.
Running the harness against the real pipeline surfaced a third candidate --- a
hash collision in the deterministic embedding. The fix was the assertion, from
an exact-rank claim to a relative-order claim, not the pipeline, which was
correct by its own contract.

\vspace{0.2cm}
\noindent\finding{F-2.} An early design for enforcing C-1 and C-2 would have
put a second, independent filter check inside the index decorator. Reasoning
through what the mutation pass would need to prove --- before any mutation
code existed --- showed this would make the intended defect structurally
unfalsifiable: a stage that built a wrong filter would still be caught by the
decorator re-checking the same filter object, so the mutation could never leak
anything, and the test proving the scenario catches the mutation would pass
for the wrong reason.

F-2 is the one to remember, because it was caught by design review rather than
by a failing test, and because the defect it prevented was not in the pipeline
at all. It was in the instrument: a measurement that would have reported
success for a reason unrelated to the thing being measured.

\subsection{The honest caveat}

The same person wrote the scenario corpus and the four broken variants, in the
same sitting, with full knowledge of how each mutation would break the
pipeline. All four were caught on the first run --- which is weaker evidence
than it looks. It proves the assertions are internally consistent with the
design intent, not that they would catch a bug nobody had in mind while
writing them. Closing that gap (D-7) needs a second, independent author, not
more code.

% ============================================================================
\section{CI and governance}

\begin{table}[htbp]
\centering
\caption{What runs, when, and what it needs.}
\begin{tabularx}{\linewidth}{@{}llX@{}}
\toprule
\textbf{Workflow} & \textbf{Trigger} & \textbf{What it does} \\
\midrule
\texttt{ci.yml} & every push and PR &
\texttt{lint-docs} (markdown, diagram parity, Ajv scenario and fixture schema
validation), \texttt{build-test} (unit tests, Layer 1, the mutation pass,
judge machinery), \texttt{secret-scan} (gitleaks over the full history) ---
zero credentials required \\
\addlinespace
\texttt{nightly.yml} & scheduled, keyed &
the full Layer 2 judged set; reports \texttt{NOT RUN} without a configured
\texttt{LLM\_API\_KEY}, honestly, rather than a silent skip nobody notices \\
\addlinespace
\texttt{build-overview-pdf.yml} & \texttt{workflow\_dispatch} only &
renders this paper and its Polish edition to PDF \\
\bottomrule
\end{tabularx}
\end{table}

There is deliberately no deploy workflow. Section~\ref{sec:readiness} explains
why, and what would need to exist first.

% ============================================================================
\section{Production readiness, honestly}
\label{sec:readiness}

This is a proof of concept, not a deployed service, and this section says
exactly what that means rather than leaving a reader to infer it from an
absent demo link.

\textbf{What is real:} the pipeline, the eval harness, the mutation pass, and
the CI that runs all three on every push, with zero credentials configured.

\textbf{What is not, yet} --- each dated and reasoned in
\texttt{docs/DEVIATIONS.md}:

\begin{table}[htbp]
\centering
\caption{The open deviations, as of this document's source. Each carries a
date and a stated condition that would close it.}
\begin{tabularx}{\linewidth}{@{}lX@{}}
\toprule
\textbf{\#} & \textbf{What has never happened} \\
\midrule
D-1 & No live embedding model or LLM judge has ever run. The vector path is a deterministic hash; the judge has never scored a real response. \\
D-2 & No corpus at production scale --- 28 scenarios against a $\sim$20-listing fixture. \\
D-3 & No real ingestion traffic. Every event was authored by a scenario file and replayed through the harness. \\
D-4 & No live Elasticsearch cluster, ever. \texttt{ElasticsearchIndex.cs} is written from the SDK's documented shapes. \\
D-5 & No production loop. No deployed, trafficked endpoint exists to extract an incident from, and no deploy workflow exists at all. \\
D-6 & The C-7 ranking-manipulation defence has been attacked only by the two scenarios planned for it, not a wider adversarial corpus. \\
D-7 & The mutation pass was not independently authored --- Section~\ref{sec:findings}'s caveat. \\
D-8 & No spelling correction or query rewriting exists. \\
\bottomrule
\end{tabularx}
\end{table}

A public roadmap issue tracks what closing each of these --- plus what a
genuinely production-ready deployment needs beyond them (authentication, a
deploy pipeline, load testing at scale, real observability) --- would take,
organised into ordered workstreams rather than a flat wishlist. Three
self-contained pieces of it have already landed without needing any external
infrastructure: rate limiting and a bounded page size on \texttt{POST /search},
the \texttt{ISearchIndex.HealthAsync()} readiness-probe wiring, and the
ingestion reordering-and-dead-letter behaviour.

% ============================================================================
\section{Decisions, deviations, and non-goals}

\subsection{Seven decisions, each naming what it rejected}

One file per decision, numbered sequentially, never renumbered or rewritten
after acceptance.

\begin{table}[htbp]
\centering
\caption{Architecture decision records (\texttt{docs/adr/}).}
\begin{tabularx}{\linewidth}{@{}llX@{}}
\toprule
\textbf{\#} & \textbf{Status} & \textbf{Title} \\
\midrule
0001 & Accepted & Record architecture decisions \\
0002 & Accepted & Mock index first, zero credentials by default \\
0003 & Accepted & Assertions read the structural trace, never response text \\
0004 & Accepted & Pin the embedding and the judge model separately, never fall back silently \\
0005 & Accepted & The Elasticsearch and vector-store SDK lives behind a five-method seam \\
0006 & Accepted & Event idempotency at the consumer boundary \\
0007 & Accepted & Render this document to PDF on demand, never committed \\
\bottomrule
\end{tabularx}
\end{table}

\subsection{Deviations, and what is deliberately not carried over}

\texttt{docs/DEVIATIONS.md} is the authoritative, dated list --- open from the
first commit, not written retroactively. Beyond the eight rows above it also
records two extensions this repository proposes back to the standards it
follows (a mutation-testing requirement for eval suites generally, and a
fixture-composition pattern for eval scenarios), and three patterns the worked
example this repository mirrors carries that this repository deliberately does
not: a confirmation-gate write boundary (there is no write path a human
approves --- the only write is an automated consumer), a showcase frontend
(this repository is about backend and ranking engineering, not UI), and a
tag-driven public deployment (no public demo exists for this proof of
concept).

\subsection{Non-goals}

Stated so that scope creep has something to fail against:

\begin{itemize}[leftmargin=1.4em]
  \item No UI. A REST API and nothing else.
  \item No personalisation, no authentication beyond a single fictional owner
        directory, no multi-tenant catalogue.
  \item No spelling correction or query rewriting (D-8).
  \item No fork of the standards this repository follows. Deviations are
        recorded, not worked around.
  \item No real-world listing data, ever, in any fixture. Every listing, owner
        and query is synthetic.
  \item No HTTP route writes to the index. The only write path is
        \texttt{IngestionConsumer}, reachable only from an ingestion event.
\end{itemize}

% ============================================================================
\section{Engineering values}

\begin{itemize}[leftmargin=1.4em]
  \item \textbf{The spec is the contract; the code is measured against it, not
        the other way round.}
  \item \textbf{A claim without a measurement is marketing.} Every number in
        this paper traces back to a file the eval suite itself produces, or to
        a document version-controlled alongside the code it describes.
  \item \textbf{Absence is worth proving, not just presence.} 28 of 128
        Layer 1 assertions (22\%) assert that something never happened.
  \item \textbf{A test suite that has never failed is a suite nobody has
        tested.} ``The constraint scenarios pass'' and ``the constraint
        scenarios would catch a violation'' are different claims, and only one
        of them is provable by watching green CI.
  \item \textbf{An honest gap, dated and reasoned, is a decision. An unstated
        one is drift.} ``We know this isn't production-ready'' is not the same
        sentence as writing down which specific parts aren't, why, and what
        would close each one.
\end{itemize}

% ============================================================================
\section{Status, and how this relates to the standards}

Complete, tested, and honestly documented as a proof of concept: the full
pipeline, 128 Layer 1 assertions across 28 scenarios, the four-variant
mutation pass, and CI that runs all of it with zero credentials configured.
Not complete, and not claimed to be, as a production service.

The forward plan lives as a public GitHub issue rather than in this document,
so it can be checked off in place rather than going stale here: eight ordered
workstreams, from a real Elasticsearch backend and a real embedding model
through a deploy pipeline, load testing at scale, and closing the two
eval-rigor gaps that need a second, independent author rather than more code.

This repository does not re-derive its architecture. It reads
\texttt{architecture-standards} and follows it --- .NET Aspire with the
AppHost as composition root, one thin \texttt{ServiceDefaults} kernel,
OpenTelemetry first --- and mirrors, rather than re-derives, the eval-bench
shape that \texttt{agent-eval-bench}, the same author's reference
implementation of the standards' AI-evaluation guide, already demonstrates for
a different specimen. \finding{The specimen changed, from a tool-using agent
to a hybrid ranker; the instrument's shape did not.} Where this repository
must depart from either, the departure is recorded, dated and reasoned rather
than implied by the absence of a pattern the worked example carries.

% ============================================================================
\section{Conclusion}

The claim this repository makes is narrow and it is measured: a hybrid ranker
respects its hard constraints on both retrieval paths, is honest about
attribution and about degradation, and does not let a listing buy its own rank
--- and each of those is bound to a named scenario that would fail if it
stopped being true. The mutation pass exists because a suite that has never
failed has not been tested, and the deviations exist because the difference
between a proof of concept and a production service is a list of specific
absent things, not a disclaimer.

What is not claimed is equally explicit. No real cluster, no trained
embedding, no judged run with a credential, no traffic, no deployment. Those
are not omissions discovered by a reader; they are rows in a dated table, each
with the condition that would close it.

\begin{thebibliography}{9}
\bibitem{spec} \texttt{docs/SPEC.md} --- the behaviour contract: twelve
behaviours, seven hard constraints, four rubrics, and the trace-event table.
Written before the pipeline.
\bibitem{findings} \texttt{docs/FINDINGS.md} --- the recomputed counts, the
mutation-pass results, and what the suite has \emph{not} caught.
\bibitem{deviations} \texttt{docs/DEVIATIONS.md} --- the dated deviations,
the proposed extensions, and the patterns deliberately not carried over.
\bibitem{overview} \texttt{docs/OVERVIEW.md} --- the Markdown source this
paper presents.
\bibitem{standards} \texttt{architecture-standards} ---
\url{https://github.com/konradcinkusz/architecture-standards}
\bibitem{aeb} \texttt{agent-eval-bench} ---
\url{https://github.com/konradcinkusz/agent-eval-bench}
\end{thebibliography}

\end{document}
